% Generated fragments require trusted human review before compilation.
% This binder never requests shell escape or runs external tools.
\documentclass[11pt]{article}
% Trusted shared apparatus. Requires LuaLaTeX; never use shell escape.
\usepackage{fontspec}
\defaultfontfeatures{Ligatures=TeX}
\IfFontExistsTF{Noto Serif}{\setmainfont{Noto Serif}}{%
  \IfFontExistsTF{DejaVu Serif}{\setmainfont{DejaVu Serif}}{%
    \IfFontExistsTF{Times New Roman}{\setmainfont{Times New Roman}}{%
      \PackageError{gelfand-binder}{No configured Cyrillic serif font installed}{Install Noto Serif or DejaVu Serif, or select an installed Cyrillic font.}%
    }%
  }%
}
\IfFontExistsTF{Noto Sans}{\setsansfont{Noto Sans}}{%
  \IfFontExistsTF{DejaVu Sans}{\setsansfont{DejaVu Sans}}{\setsansfont{Arial}}%
}
\IfFontExistsTF{Noto Sans Mono}{\setmonofont{Noto Sans Mono}}{%
  \IfFontExistsTF{DejaVu Sans Mono}{\setmonofont{DejaVu Sans Mono}}{\setmonofont{Courier New}}%
}
\usepackage[a4paper,margin=22mm]{geometry}
\usepackage{amsmath,amssymb,amsthm,mathtools}
\usepackage{xcolor}
\usepackage{tikz}
\usepackage{tikz-cd}
\usetikzlibrary{arrows.meta,positioning,calc}
\usepackage{hyperref}
\hypersetup{unicode=true,colorlinks=true,linkcolor=blue,pdftitle={Gelfand transcription binder}}
\newtheorem{theorem}{Theorem}
\newtheorem{lemma}{Lemma}
\newtheorem{proposition}{Proposition}
\newtheorem{corollary}{Corollary}
\theoremstyle{definition}
\newtheorem{definition}{Definition}
\newtheorem{example}{Example}
\theoremstyle{remark}
\newtheorem{remark}{Remark}
\setlength{\parindent}{0pt}
\setlength{\parskip}{0.55em}
\setlength{\emergencystretch}{3em}
\newcommand{\SourcePage}[2]{%
  \clearpage\phantomsection\label{#1:source-page}%
  \addcontentsline{toc}{subsection}{#1}%
  {\small\sffamily\color{blue!60!black}Source scan frame: #1; visible page/leaf label: #2}\par\medskip%
}
\newcommand{\Unclear}[3]{%
  \begingroup\color{orange!70!black}%
  [Unresolved reading #1/#2: #3]\endgroup%
}
\newcommand{\Incomplete}[3]{%
  \par\begingroup\color{red!65!black}\textbf{Incomplete #1/#2:} #3\endgroup\par%
}
\newcommand{\SourceFigure}[3]{%
  \par\begingroup\color{blue!60!black}\textbf{\hyperref[#1:source-page]{Source figure #1/#2}:} #3%
  \par\textit{Consult the authoritative scan at this locator.}\endgroup\par%
}



% Reader layout only. Original diplomatic text is retained verbatim.
\setlength{\parskip}{0.35em}
\setlength{\emergencystretch}{4em}
\renewcommand{\SourcePage}[2]{%
  \par\addvspace{1em}\phantomsection\label{#1:source-page}%
  \addcontentsline{toc}{subsection}{#1}%
  {\footnotesize\sffamily\color{blue!55!black}#1%
    \def\givenlabel{#2}\def\nolabel{unlabelled}%
    \ifx\givenlabel\nolabel\else\quad[номер в источнике: #2]\fi
  }\par\nobreak\smallskip%
}
\renewcommand{\Unclear}[3]{%
  \begingroup\color{orange!65!black}%
  \ifmmode\text{[неясно: #3]}\else[неясно: #3]\fi\endgroup%
}
\renewcommand{\Incomplete}[3]{%
  \par\begingroup\color{red!60!black}\textbf{Неразобранный фрагмент:} #3%
  \endgroup\par%
}
\renewcommand{\SourceFigure}[3]{%
  \par\begingroup\color{blue!55!black}\textbf{Рисунок:} #3%
  \par{\small\textit{Исходный рисунок: \hyperref[#1:source-page]{#1}, #2.}}%
  \endgroup\par%
}
\hypersetup{pdftitle={Семинар И. М. Гельфанда. Дипломатическая транскрипция},pdfauthor={},pdfsubject={Русское рабочее издание с отмеченными сомнительными чтениями}}

\usepackage[normalem]{ulem}
\usepackage{newunicodechar}
\usepackage{adjustbox}
\newfontfamily\SourceSymbolFont{Segoe UI Symbol}
\newfontfamily\SourceCJKFont{MS Gothic}
\let\DiplomaticMathUnderline\underline
\renewcommand{\underline}[1]{\ifmmode\DiplomaticMathUnderline{#1}\else\uline{#1}\fi}
\newunicodechar{→}{\ifmmode\text{{\SourceSymbolFont\char"2192}}\else{\SourceSymbolFont\char"2192}\fi}
\newunicodechar{⇒}{\ifmmode\text{{\SourceSymbolFont\char"21D2}}\else{\SourceSymbolFont\char"21D2}\fi}
\newunicodechar{∀}{\ifmmode\text{{\SourceSymbolFont\char"2200}}\else{\SourceSymbolFont\char"2200}\fi}
\newunicodechar{∂}{\ifmmode\text{{\SourceSymbolFont\char"2202}}\else{\SourceSymbolFont\char"2202}\fi}
\newunicodechar{∈}{\ifmmode\text{{\SourceSymbolFont\char"2208}}\else{\SourceSymbolFont\char"2208}\fi}
\newunicodechar{∌}{\ifmmode\text{{\SourceSymbolFont\char"220C}}\else{\SourceSymbolFont\char"220C}\fi}
\newunicodechar{√}{\ifmmode\text{{\SourceSymbolFont\char"221A}}\else{\SourceSymbolFont\char"221A}\fi}
\newunicodechar{∞}{\ifmmode\text{{\SourceSymbolFont\char"221E}}\else{\SourceSymbolFont\char"221E}\fi}
\newunicodechar{∧}{\ifmmode\text{{\SourceSymbolFont\char"2227}}\else{\SourceSymbolFont\char"2227}\fi}
\newunicodechar{∩}{\ifmmode\text{{\SourceSymbolFont\char"2229}}\else{\SourceSymbolFont\char"2229}\fi}
\newunicodechar{∪}{\ifmmode\text{{\SourceSymbolFont\char"222A}}\else{\SourceSymbolFont\char"222A}\fi}
\newunicodechar{∫}{\ifmmode\text{{\SourceSymbolFont\char"222B}}\else{\SourceSymbolFont\char"222B}\fi}
\newunicodechar{∴}{\ifmmode\text{{\SourceSymbolFont\char"2234}}\else{\SourceSymbolFont\char"2234}\fi}
\newunicodechar{∼}{\ifmmode\text{{\SourceSymbolFont\char"223C}}\else{\SourceSymbolFont\char"223C}\fi}
\newunicodechar{≈}{\ifmmode\text{{\SourceSymbolFont\char"2248}}\else{\SourceSymbolFont\char"2248}\fi}
\newunicodechar{≠}{\ifmmode\text{{\SourceSymbolFont\char"2260}}\else{\SourceSymbolFont\char"2260}\fi}
\newunicodechar{≤}{\ifmmode\text{{\SourceSymbolFont\char"2264}}\else{\SourceSymbolFont\char"2264}\fi}
\newunicodechar{≥}{\ifmmode\text{{\SourceSymbolFont\char"2265}}\else{\SourceSymbolFont\char"2265}\fi}
\newunicodechar{⊕}{\ifmmode\text{{\SourceSymbolFont\char"2295}}\else{\SourceSymbolFont\char"2295}\fi}
\newunicodechar{⊗}{\ifmmode\text{{\SourceSymbolFont\char"2297}}\else{\SourceSymbolFont\char"2297}\fi}
\newunicodechar{①}{\ifmmode\text{{\SourceSymbolFont\char"2460}}\else{\SourceSymbolFont\char"2460}\fi}
\newunicodechar{②}{\ifmmode\text{{\SourceSymbolFont\char"2461}}\else{\SourceSymbolFont\char"2461}\fi}
\newunicodechar{③}{\ifmmode\text{{\SourceSymbolFont\char"2462}}\else{\SourceSymbolFont\char"2462}\fi}
\newunicodechar{④}{\ifmmode\text{{\SourceSymbolFont\char"2463}}\else{\SourceSymbolFont\char"2463}\fi}
\newunicodechar{Ⓗ}{\ifmmode\text{{\SourceSymbolFont\char"24BD}}\else{\SourceSymbolFont\char"24BD}\fi}
\newunicodechar{〃}{\ifmmode\text{{\SourceCJKFont\char"3003}}\else{\SourceCJKFont\char"3003}\fi}


% Reader revision 3. Annotation wording and candidates remain unchanged.
\usetikzlibrary{shapes.geometric,decorations.pathreplacing,decorations.markings,patterns,fit}
\newsavebox{\ReaderAnnotationBox}
\newcommand{\ReaderAnnotation}[1]{%
  \sbox{\ReaderAnnotationBox}{[неясно: #1]}%
  \ifdim\wd\ReaderAnnotationBox>.70\linewidth
    \parbox[t]{.70\linewidth}{\raggedright[неясно: #1]}%
  \else\usebox{\ReaderAnnotationBox}\fi%
}
\ExplSyntaxOn
\cs_new_protected:Npn \reader_unclear_candidate:n #1
 {
  \tl_set:Nx \l_tmpa_tl { \tl_to_str:n {#1} }
  % Text explicitly guarded inside a mathematical candidate remains math-safe.
  \regex_replace_all:nnN { \\(?:text|textrm|textnormal|textbf|operatorname)\s*\{[^{}]*\} } { } \l_tmpa_tl
  \regex_match:nVTF { [\x{0400}-\x{052F}] | \$ } \l_tmpa_tl
   { \ifmmode \text{\ReaderAnnotation{#1}} \else \ReaderAnnotation{#1} \fi }
   {
    \ifmmode
      [\text{неясно:~}#1]
    \else
      \regex_match:nVTF { [_^] | \\ } \l_tmpa_tl
       { \ensuremath{[\text{неясно:~}#1]} }
       { \ReaderAnnotation{#1} }
    \fi
   }
 }
\RenewDocumentCommand \Unclear { m m m }
 { \begingroup\color{orange!65!black}\reader_unclear_candidate:n{#3}\endgroup }
\ExplSyntaxOff
\newsavebox{\ReaderIncompleteBox}
\newcommand{\ReaderIncompleteNote}[1]{%
 \sbox{\ReaderIncompleteBox}{\textbf{Неразобранный фрагмент:} #1}%
 \ifdim\wd\ReaderIncompleteBox>.70\linewidth
  \parbox[t]{.70\linewidth}{\raggedright\textbf{Неразобранный фрагмент:} #1}%
 \else\usebox{\ReaderIncompleteBox}\fi%
}
\renewcommand{\Incomplete}[3]{%
 \ifmmode\text{\color{red!60!black}\ReaderIncompleteNote{#3}}%
 \else\par\begingroup\color{red!60!black}\textbf{Неразобранный фрагмент:} #3\endgroup\par\fi%
}
% Ragged right avoids stretching long untranslated abbreviations into margins.
\AtBeginDocument{\raggedright}
\newunicodechar{А}{\ifmmode\text{\char"410\relax}\else\char"410\relax\fi}
\newunicodechar{Б}{\ifmmode\text{\char"411\relax}\else\char"411\relax\fi}
\newunicodechar{В}{\ifmmode\text{\char"412\relax}\else\char"412\relax\fi}
\newunicodechar{Г}{\ifmmode\text{\char"413\relax}\else\char"413\relax\fi}
\newunicodechar{Д}{\ifmmode\text{\char"414\relax}\else\char"414\relax\fi}
\newunicodechar{Е}{\ifmmode\text{\char"415\relax}\else\char"415\relax\fi}
\newunicodechar{Ж}{\ifmmode\text{\char"416\relax}\else\char"416\relax\fi}
\newunicodechar{З}{\ifmmode\text{\char"417\relax}\else\char"417\relax\fi}
\newunicodechar{И}{\ifmmode\text{\char"418\relax}\else\char"418\relax\fi}
\newunicodechar{Й}{\ifmmode\text{\char"419\relax}\else\char"419\relax\fi}
\newunicodechar{К}{\ifmmode\text{\char"41A\relax}\else\char"41A\relax\fi}
\newunicodechar{Л}{\ifmmode\text{\char"41B\relax}\else\char"41B\relax\fi}
\newunicodechar{М}{\ifmmode\text{\char"41C\relax}\else\char"41C\relax\fi}
\newunicodechar{Н}{\ifmmode\text{\char"41D\relax}\else\char"41D\relax\fi}
\newunicodechar{О}{\ifmmode\text{\char"41E\relax}\else\char"41E\relax\fi}
\newunicodechar{П}{\ifmmode\text{\char"41F\relax}\else\char"41F\relax\fi}
\newunicodechar{Р}{\ifmmode\text{\char"420\relax}\else\char"420\relax\fi}
\newunicodechar{С}{\ifmmode\text{\char"421\relax}\else\char"421\relax\fi}
\newunicodechar{Т}{\ifmmode\text{\char"422\relax}\else\char"422\relax\fi}
\newunicodechar{У}{\ifmmode\text{\char"423\relax}\else\char"423\relax\fi}
\newunicodechar{Ф}{\ifmmode\text{\char"424\relax}\else\char"424\relax\fi}
\newunicodechar{Х}{\ifmmode\text{\char"425\relax}\else\char"425\relax\fi}
\newunicodechar{Ц}{\ifmmode\text{\char"426\relax}\else\char"426\relax\fi}
\newunicodechar{Ч}{\ifmmode\text{\char"427\relax}\else\char"427\relax\fi}
\newunicodechar{Ш}{\ifmmode\text{\char"428\relax}\else\char"428\relax\fi}
\newunicodechar{Щ}{\ifmmode\text{\char"429\relax}\else\char"429\relax\fi}
\newunicodechar{Ъ}{\ifmmode\text{\char"42A\relax}\else\char"42A\relax\fi}
\newunicodechar{Ы}{\ifmmode\text{\char"42B\relax}\else\char"42B\relax\fi}
\newunicodechar{Ь}{\ifmmode\text{\char"42C\relax}\else\char"42C\relax\fi}
\newunicodechar{Э}{\ifmmode\text{\char"42D\relax}\else\char"42D\relax\fi}
\newunicodechar{Ю}{\ifmmode\text{\char"42E\relax}\else\char"42E\relax\fi}
\newunicodechar{Я}{\ifmmode\text{\char"42F\relax}\else\char"42F\relax\fi}
\newunicodechar{а}{\ifmmode\text{\char"430\relax}\else\char"430\relax\fi}
\newunicodechar{б}{\ifmmode\text{\char"431\relax}\else\char"431\relax\fi}
\newunicodechar{в}{\ifmmode\text{\char"432\relax}\else\char"432\relax\fi}
\newunicodechar{г}{\ifmmode\text{\char"433\relax}\else\char"433\relax\fi}
\newunicodechar{д}{\ifmmode\text{\char"434\relax}\else\char"434\relax\fi}
\newunicodechar{е}{\ifmmode\text{\char"435\relax}\else\char"435\relax\fi}
\newunicodechar{ж}{\ifmmode\text{\char"436\relax}\else\char"436\relax\fi}
\newunicodechar{з}{\ifmmode\text{\char"437\relax}\else\char"437\relax\fi}
\newunicodechar{и}{\ifmmode\text{\char"438\relax}\else\char"438\relax\fi}
\newunicodechar{й}{\ifmmode\text{\char"439\relax}\else\char"439\relax\fi}
\newunicodechar{к}{\ifmmode\text{\char"43A\relax}\else\char"43A\relax\fi}
\newunicodechar{л}{\ifmmode\text{\char"43B\relax}\else\char"43B\relax\fi}
\newunicodechar{м}{\ifmmode\text{\char"43C\relax}\else\char"43C\relax\fi}
\newunicodechar{н}{\ifmmode\text{\char"43D\relax}\else\char"43D\relax\fi}
\newunicodechar{о}{\ifmmode\text{\char"43E\relax}\else\char"43E\relax\fi}
\newunicodechar{п}{\ifmmode\text{\char"43F\relax}\else\char"43F\relax\fi}
\newunicodechar{р}{\ifmmode\text{\char"440\relax}\else\char"440\relax\fi}
\newunicodechar{с}{\ifmmode\text{\char"441\relax}\else\char"441\relax\fi}
\newunicodechar{т}{\ifmmode\text{\char"442\relax}\else\char"442\relax\fi}
\newunicodechar{у}{\ifmmode\text{\char"443\relax}\else\char"443\relax\fi}
\newunicodechar{ф}{\ifmmode\text{\char"444\relax}\else\char"444\relax\fi}
\newunicodechar{х}{\ifmmode\text{\char"445\relax}\else\char"445\relax\fi}
\newunicodechar{ц}{\ifmmode\text{\char"446\relax}\else\char"446\relax\fi}
\newunicodechar{ч}{\ifmmode\text{\char"447\relax}\else\char"447\relax\fi}
\newunicodechar{ш}{\ifmmode\text{\char"448\relax}\else\char"448\relax\fi}
\newunicodechar{щ}{\ifmmode\text{\char"449\relax}\else\char"449\relax\fi}
\newunicodechar{ъ}{\ifmmode\text{\char"44A\relax}\else\char"44A\relax\fi}
\newunicodechar{ы}{\ifmmode\text{\char"44B\relax}\else\char"44B\relax\fi}
\newunicodechar{ь}{\ifmmode\text{\char"44C\relax}\else\char"44C\relax\fi}
\newunicodechar{э}{\ifmmode\text{\char"44D\relax}\else\char"44D\relax\fi}
\newunicodechar{ю}{\ifmmode\text{\char"44E\relax}\else\char"44E\relax\fi}
\newunicodechar{я}{\ifmmode\text{\char"44F\relax}\else\char"44F\relax\fi}
\newunicodechar{Ё}{\ifmmode\text{\char"401\relax}\else\char"401\relax\fi}
\newunicodechar{ё}{\ifmmode\text{\char"451\relax}\else\char"451\relax\fi}

\begin{document}
\begin{center}
{\Large Семинар И. М. Гельфанда}\par\medskip
{\LARGE Дипломатическая транскрипция}\par\bigskip
{\large Рабочее русское издание}\par\medskip
{\large Том 13 из 34. Архивная глава 013}\par\bigskip
Кадры GELFAND-P000777 -- GELFAND-P000785\par

\end{center}
\medskip
Сокращения, отменённые записи и сомнительные чтения сохранены.
Неразобранные фрагменты отмечены в тексте. Математические утверждения
источника переданы без молчаливых исправлений.

Это рабочее издание: сверка с рукописью продолжается. В текущем комплекте
учтены 1503 из 1527 кадров; 24 кадра пока представлены явными указателями
отсутствующего возврата. Полный журнал чтений и происхождения данных
сохранён в основном рабочем переплёте.

Границы кадров показаны небольшими синими указателями; печатный текст
продолжается без обязательного разрыва страницы после каждого кадра.
\bigskip\hrule\bigskip
\SourcePage{GELFAND-P000777}{unlabelled}
% Additional independent review: original image directly inspected.
% Enlargement resamples the original pixels; it does not add source detail.
% REVIEW independent eq1-overwrite: box=[1000, 900, 1210, 1030], scale=6x, interpolation=nearest; original-pixel coordinates.
% REVIEW original directly viewed: 2408x3448 pixels. Coordinate origin top-left; crop boxes [x0,y0,x1,y1), right/bottom exclusive; PIL LANCZOS enlargement.
% REVIEW heading: box=[185, 550, 1450, 880], scale=3x.
% REVIEW eq1-3: box=[140, 885, 2110, 1470], scale=3x.
% REVIEW scattering: box=[340, 1490, 2340, 2150], scale=3x.
% REVIEW lower-prose: box=[280, 2160, 2260, 2860], scale=3x.
% REVIEW eq5: box=[115, 2890, 2408, 3100], scale=3x.
% REVIEW eq1-sign: box=[1045, 925, 1138, 1018], scale=10x.
% REVIEW gutter: box=[0, 2500, 87, 2690], scale=3x.
% REVIEW pole-rho: box=[808, 1980, 1040, 2070], scale=10x.
% REVIEW eq6-psi-argument: box=[1630, 3110, 1810, 3250], scale=10x.
% REVIEW cancelled-line: box=[325, 2170, 1070, 2270], scale=3x.
% REVIEW cancelled-variable: box=[580, 2280, 770, 2400], scale=6x.
% REVIEW cancelled-heading: box=[600, 565, 890, 645], scale=6x.
% REVIEW technique: box=[305, 2635, 2180, 2750], scale=3x.
% REVIEW gutter-context: box=[0, 2480, 120, 2680], scale=5x.
\underline{27/III-78.}
[Зачёркнуто: \Incomplete{GELFAND-P000777}{heading-cancelled}{Текст между датой и словами «Сем. Гельфанда» густо зачёркнут; читается только начальная С.}]
\underline{Сем. Гельфанда.}

\underline{Ф. Калоджеро.}\qquad \underline{О преобр. Бэклунда.}

(всё это вместе с A. Degasperis)

\begin{equation}
\Psi_{xx}(x,k)=\bigl[q(x)\,\Unclear{GELFAND-P000777}{eq01-sign}{\text{неразборчивый знак}}\,k^2\bigr]\Psi(x,k),\qquad -\infty<x<+\infty
\tag{1}
\end{equation}
\begin{equation}
q(x)\to0\quad\text{при }|x|\to\infty\quad\text{дост. быстро}
\tag{2}
\end{equation}
Пусть $\Psi$ — такое реш., что
\begin{equation}
\left\{\begin{aligned}
\Psi(x,k)&\underset{x\to+\infty}{\longrightarrow}e^{-ikx}+R(k)e^{ikx},\\
\Psi(x,k)&\underset{x\to-\infty}{\longrightarrow}T(k)e^{-ikx}
\end{aligned}\right.
\tag{3}
\end{equation}

\begin{equation*}
q(x)\longmapsto S:\{R(k),\ -\infty<k<+\infty;\ \text{и нек-рая инф. о дискр. спектре}\}.
\end{equation*}
[Под $S$, с соединительной чертой: спектр. преобр.]

Для широкого класса потенциалов дискр. сп. = полюса $R(k)$:
\begin{equation}
R(k)\sim\frac{ic_n^2}{k-i\rho_n}
\tag{4}
\end{equation}
Потенциал $q$ определяется по $R(k)$.

[Зачёркнуто: Сущ. вз. однозн. соотв.]

Пусть [зачёркнуто: \Unclear{GELFAND-P000777}{cancelled-before-q}{$q(x)$}] $q=q(x,t)$, $R=R(k,t)$.

\underline{В каких случаях}

$R$ мен. просто, а $q$ мен. интересно?

[У левого края кадра, за линией поля: \Unclear{GELFAND-P000777}{gutter-fragment}{…а))}]

\underline{Техника:} т.к. $R(k)$ опр. асимпт., то естеств. исп. \underline{вронскианы}.

\underline{2 ур-я.} Пусть есть $q,q'\longleftrightarrow R,R'$.

\begin{equation}
2ikf(-4k^2)[R'(k)-R(k)]=\int_{-\infty}^{\infty}dx\,\Psi'(x,k)\Psi(x,k)f(\underline{\Lambda})[q'(x)-q(x)]
\tag{5}
\end{equation}
\begin{equation}
(2ik)^2g(-4k^2)[R'(k)+R(k)]=\int_{-\infty}^{\infty}dx\,\Psi'(x,k)\Unclear{GELFAND-P000777}{eq06-psi-argument}{\Psi(x,t)}g(\underline{\Lambda})\underline{\Gamma}\cdot1
\tag{6}
\end{equation}

\SourcePage{GELFAND-P000778}{unlabelled}
% Additional independent review: original image directly inspected.
% Enlargement resamples the original pixels; it does not add source detail.
% REVIEW independent top-operator-dependence: box=[910, 50, 1530, 180], scale=3x, interpolation=nearest; original-pixel coordinates.
% REVIEW independent bottom-line: box=[310, 3320, 1450, 3448], scale=3x, interpolation=nearest; original-pixel coordinates.
% REVIEW independent alternate-operator: box=[300, 3090, 880, 3310], scale=3x, interpolation=nearest; original-pixel coordinates.
% REVIEW independent alternate-operator-right: box=[840, 3090, 1420, 3310], scale=3x, interpolation=nearest; original-pixel coordinates.
% REVIEW original directly viewed: 2408x3448 pixels. Coordinate origin top-left; crop boxes [x0,y0,x1,y1), right/bottom exclusive; PIL LANCZOS enlargement.
% REVIEW top-definitions: box=[180, 45, 2040, 375], scale=3x.
% REVIEW eq7-8: box=[65, 355, 2360, 755], scale=3x.
% REVIEW upper-middle: box=[200, 740, 2270, 1200], scale=3x.
% REVIEW middle: box=[100, 1210, 1930, 1950], scale=3x.
% REVIEW lower-prose: box=[250, 1940, 1730, 2840], scale=3x.
% REVIEW lower-formulas: box=[90, 2800, 2200, 3448], scale=3x.
% REVIEW top-variables: box=[1280, 54, 1530, 153], scale=10x.
% REVIEW lower-cancellation: box=[300, 3260, 1430, 3448], scale=3x.
% REVIEW bottom-words: box=[825, 3310, 1440, 3442], scale=6x.
% REVIEW alternative-L: box=[360, 3130, 1390, 3300], scale=5x.
% REVIEW small-cancel-property: box=[259, 838, 400, 933], scale=6x.
% REVIEW small-cancel-parens: box=[341, 1815, 450, 1920], scale=6x.
% REVIEW small-cancel-gamma: box=[565, 2700, 657, 2795], scale=6x.
% REVIEW first-definition-word: box=[388, 76, 854, 166], scale=6x.
% REVIEW variables-context: box=[1170, 50, 1535, 180], scale=5x.
$\underline{\Lambda}$ — интегро-диф. опр., зав. от \Unclear{GELFAND-P000778}{top-variables}{$u$ и $u'$}

$\underline{\Gamma}$ — тоже оператор

Здесь
\begin{equation}
\underline{\Lambda}F(x)=F_{xx}(x)-2[q'(x)+q(x)]F(x)+\underline{\Gamma}\int_x^{\infty}dx'\,F(x');
\tag{7}
\end{equation}
\begin{equation}
\underline{\Gamma}G(x)=[q'_x(x)+q_x(x)]G(x)+[q'(x)-q(x)]\int_x^{\infty}dx'[q'(x')-q(x')]G(x');
\tag{8}
\end{equation}
$f,g$ — произв. ф-и.

[Зачёркнуто: \Unclear{GELFAND-P000778}{cancelled-before-property}{неразборчиво}]
Свойство $\underline{\Lambda}$:
\begin{equation}
\int_{-\infty}^{\infty}dx\,\Psi'(x,k)\Psi(x,k)\underline{\Lambda}F(x)=-4k^2\int_{-\infty}^{+\infty}dx\,\Psi'(x,k)\Psi(x,k)F(x).
\tag{9}
\end{equation}
Пусть $q$ и $q'$ различны (потом будет взят предел).

Если $q$ и $q'$ таковы, что
\begin{equation}
f(\underline{\Lambda})[q'(x)-q(x)]+g(\underline{\Lambda})\underline{\Gamma}\,1=0,
\tag{10}
\end{equation}
то
\begin{equation}
R'(k)=\frac{f(-4k^2)-2ikg(-4k^2)}{f(-4k^2)+2ikg(-4k^2)}R(k)
\tag{11}
\end{equation}
[Зачёркнуто: \Unclear{GELFAND-P000778}{cancelled-before-parenthesis}{неразборчиво}]
(это получается из формул (5), (6))

Всё это для стац. ур-я Шрё.

Теперь надо инфинитезим. деформации брать
\begin{equation*}
q=q(x,t),\qquad q'=q(x,t+dt)
\end{equation*}
\begin{align*}
R'-R&\longmapsto R_t(k,t)\\
q'-q&\longmapsto q_t(x,t)\\
\underline{\Lambda}&\longmapsto\underline{L}\qquad(=\underline{\Lambda}|_{q'=q})
\end{align*}
$\underline{\Gamma}$ [зачёркнуто: \Unclear{GELFAND-P000778}{cancelled-before-gamma-arrow}{и}] $\longmapsto$ умн. на $2q_x(x)$

Получаем:
\begin{equation}
\underline{L}F(x)=F_{xx}(x)-4q(x,t)F(x)+2q_x(x,t)\int_x^{\infty}dx'\,F(x')
\tag{12}
\end{equation}
\underline{Гельф.} (и Дикий):
\begin{equation*}
L=\left(\frac{d^3}{dx^3}+2q\frac{d}{dx}+2\frac{d}{dx}q\right)\left(\frac{d}{dx}\right)^{-1}
\end{equation*}
\Unclear{GELFAND-P000778}{alternative-l-signs}{[В этой записи оба знака перед членами с $q$ — «+».]}

[Зачёркнуто: = опр., к-рый]

— опр., к-рый был в прошлый раз.

\SourcePage{GELFAND-P000779}{5}
% Additional independent review: original image directly inspected.
% Enlargement resamples the original pixels; it does not add source detail.
% REVIEW independent top-names-left: box=[140, 205, 1060, 330], scale=3x, interpolation=nearest; original-pixel coordinates.
% REVIEW independent top-names-right: box=[1020, 205, 1640, 330], scale=3x, interpolation=nearest; original-pixel coordinates.
% REVIEW original directly viewed: 2408x3448 pixels. Coordinate origin top-left; crop boxes [x0,y0,x1,y1), right/bottom exclusive; PIL LANCZOS enlargement.
% REVIEW top: box=[130, 28, 2200, 390], scale=3x.
% REVIEW eq13-14: box=[120, 380, 2230, 1020], scale=3x.
% REVIEW middle-prose: box=[230, 1030, 2210, 1920], scale=3x.
% REVIEW eq15-prose: box=[160, 1920, 2380, 2335], scale=3x.
% REVIEW eq16-17: box=[150, 2300, 2400, 2860], scale=3x.
% REVIEW bottom: box=[150, 2770, 2400, 3448], scale=3x.
% REVIEW scriptL: box=[349, 1035, 508, 1130], scale=10x.
% REVIEW names: box=[180, 213, 1500, 315], scale=6x.
% REVIEW diagram18: box=[380, 3160, 1240, 3430], scale=3x.
% REVIEW diagram18-rho1: box=[570, 3160, 657, 3240], scale=10x.
% REVIEW cancelled-constant: box=[1160, 1670, 1640, 1790], scale=5x.
% REVIEW inserted-eigenvalue: box=[1140, 2050, 1640, 2230], scale=5x.
% REVIEW last-prose: box=[310, 2990, 2408, 3145], scale=5x.
% REVIEW diagram-number: box=[185, 3205, 325, 3335], scale=10x.
% REVIEW top-abbrev: box=[485, 136, 675, 236], scale=10x.
% REVIEW top-satisfy: box=[615, 145, 721, 234], scale=10x.
% REVIEW small-cancel-top: box=[337, 145, 422, 229], scale=10x.
% REVIEW dots-row: box=[690, 728, 1340, 795], scale=5x.
% REVIEW cancelled-eigenword: box=[1330, 2121, 1556, 2195], scale=10x.
% REVIEW surname-tight: box=[1080, 218, 1390, 310], scale=6x.
% REVIEW fL-line: box=[485, 755, 1085, 880], scale=6x.
% REVIEW eq14-polynomial-left: box=[585, 897, 840, 1015], scale=6x.
% REVIEW eq14-polynomial-right: box=[1640, 888, 1990, 1015], scale=6x.
появл. во многих местах.

а) [Зачёркнуто: \Unclear{GELFAND-P000779}{cancelled-top-symbol}{и}]
$\Psi_1,\Psi_2$ уд. ур-ю 2-го пор. $\Rightarrow\Psi_1\Psi_2$ удовл. ур-ю 3-го пор.

Гельф. Дикий; Лакс; \underline{Марченко}; \underline{\Unclear{GELFAND-P000779}{top-surname}{Мосейка}}; и др.

Св-во $\underline{\Lambda}$.
\begin{equation}
\int_{-\infty}^{\infty}dx\,\Psi'(x,k)\Psi(x,k)(\underline{\Lambda}+4k^2)F(x)=0
\tag{13}
\end{equation}
$\Psi=\Psi'\Rightarrow\underline{\Lambda}$ заменяется на $\underline{L}$.

[Над следующей строкой: ряд точек.]

Для $f(\underline{L})=\underline{L}$ имеем
\begin{equation}
q_t(x,t)=\Unclear{GELFAND-P000779}{polynomial-symbol}{\alpha}(\underline{L})q_x(x,t)\Longleftrightarrow R_t(k,t)=2ik\Unclear{GELFAND-P000779}{polynomial-symbol}{\alpha}(-4k^2)R(k,t)
\tag{14}
\end{equation}
$\Unclear{GELFAND-P000779}{polynomial-symbol}{\alpha}$ — полином $\Rightarrow KdV$ и высшие $KdV$

$R'\leftrightarrow R$ по ф-ле (11) и $R$ — реш. лин. ур-я $(14)_2$, то $R'$ тоже решение.

Тогда $q'$, постр. по $q$, удовл. ур-ю $(14)_1$ вместе с $q$.

$(14)_1$ — высшие $KdV$ в форме $AKNS$

Самое простое: пусть [зачёркнуто: $f,g=\mathrm{const}$] $g=1$, $f=2\rho$.

Получ. самое простое преобразование:
\begin{equation}
R'(k,t)=-\frac{k+i\rho}{k-i\rho}R(k,t)
\tag{15}
\end{equation}
Это добавл. солитона (доб. [зачёркнуто: \Unclear{GELFAND-P000779}{cancelled-eigenword}{неразборчиво}; вписано сверху с соединительной линией: соб. знач.] и умн. на коэффициент)

Соотн. между потенц. выписывается:

пусть
\begin{equation}
Q(x,t)=\int_x^{\infty}dx'\,q(x',t)
\tag{16}
\end{equation}
\begin{equation}
Q'_x(x,t)+Q_x(x,t)=-\frac12[Q'(x,t)-Q(x,t)][4\rho+Q'(x,t)-Q(x,t)]
\tag{17}
\end{equation}
$Q\equiv0\Rightarrow Q'$ — односолитонное решение, (для $q_t=q_{xxx}-6q_xq$)

(17) — ур-е \underline{Риккати} (отн. $Q'$ или отн. $Q$)

Все интегр. получаются послед. прим. преобр. Бэклунда из тривиального $\int u^2\,dx$

\underline{Другое применение}

\begin{equation}
\begin{adjustbox}{max width=\linewidth}
\begin{tikzcd}
Q \arrow[r,"\rho_1"] \arrow[dr,"\rho_2"'] & Q_1 \arrow[r,"\rho_2"] & Q'' \\
& Q_2 \arrow[ur,"\rho_1"'] &
\end{tikzcd}
\end{adjustbox}
\tag{18}
\end{equation}
\underline{коммутативность!}

\SourcePage{GELFAND-P000780}{unlabelled}
% Additional independent review: original image directly inspected.
% Enlargement resamples the original pixels; it does not add source detail.
% REVIEW independent superposition-sign: box=[890, 560, 1120, 680], scale=8x, interpolation=nearest; original-pixel coordinates.
Это сразу видно; т. к. на ф-ях $R$ просто умножение.
На ф-ях $q$ же нетривиально.
Получ. принцип нелин. суперпозиции:
\begin{equation*}
Q''(x,t)=Q(x,t)-
\frac{(\rho_1+\rho_2)[Q_1(x,t)-Q_2(x,t)]}
{\left\{\rho_1-\rho_2+\frac{1}{2}
[Q_1(x,t)\,\text{[знак неразборчив]}\,Q_2(x,t)]\right\}}
\tag{19}
\end{equation*}
[В рукописи дробь в (19) записана через косую черту в конце первой строки и в начале второй.]
\Unclear{GELFAND-P000780}{equation19-denominator-sign}{Между $Q_1(x,t)$ и $Q_2(x,t)$ знак густо переписан; возможны $+$ и $-$.}

\begin{center}
— · —
\end{center}

$f$ и $g$ могут зависеть от $t$!
Соотн. между (10) и (11) всё равно верно.
\begin{equation*}
q_t(x,t)=\alpha(\underline{L})q_x(x,t)
\tag{20}
\end{equation*}
\begin{equation*}
R(k,t)=e^{2ik\alpha(-4k^2)(t-t_0)}R(k,t_0)
\tag{21}
\end{equation*}
\Unclear{GELFAND-P000780}{equations20-21-polynomial-glyph}{В (20) и (21) условно передано $\alpha$; почерк допускает также $\mathcal{L}$.}

Выберем $f$ и $g$ так, чтобы дробь в (11) совпад. с выписанной
экспонентой в (21). Тогда получим соответствие
$q(x,t)$ и $q(x,t_0)$ (ф-ла (10)).

Эволюция по $t$ — также преобр. Бэкл.!

Есть аналогия с преобр. Фурье:
\begin{equation*}
u(x,t)=e^{(t-t_0)\psi(\frac{\partial}{\partial x})}u(x,t_0),
\end{equation*}
\Unclear{GELFAND-P000780}{fourier-overwritten-u}{Первое $u$ переписано поверх более раннего неразборчивого знака.}
при
\begin{equation*}
u_t(x,t)=\psi\left(\frac{\partial}{\partial x}\right)u(x,t).
\end{equation*}
\Unclear{GELFAND-P000780}{fourier-function-glyph}{В обеих строках условно передано $\psi$; возможное другое чтение — $\varphi$.}

\begin{center}
— · —
\end{center}

Рассм. ур-е
\begin{equation*}
\Psi_{tt}+\frac{1}{2}\Psi_{xx}-x\Psi_x+n\Psi=0,
\tag{22}
\end{equation*}
где $n\in\mathbb{Z}$

[Зачёркнуто, предположительно: Его. Над строкой вписано: Рассм.]
\Unclear{GELFAND-P000780}{cancelled-word-before-solution}{Его (возможно Её); зачёркивание закрывает часть букв.}
решение (вида):
\begin{equation*}
\Psi(x,t)=\prod_{j=1}^{n}(x-x_j(t))
\tag{23}
\end{equation*}
[Справа от (23) зачёркнуто: $=2^{-n}\sum_{m=0}^{n}c_m(t)H_{n-m}(x)$.]

Тогда
\begin{equation*}
\ddot{x}_j=-x_j(t)+\sum_{k=1}^{n}{}'
\frac{1+2\dot{x}_j(t)\dot{x}_k(t)}{x_j(t)-x_k(t)}
\tag{24}
\end{equation*}
(Мозер)
[Справа зачёркнуто слово; различимо лишь возможное начальное М или Н.]
\Incomplete{GELFAND-P000780}{bottom-cancelled-name}{Зачёркнутое слово справа от «Мозер» не удаётся прочесть; сохранены различимые фрагменты.}

[Ниже вписано одно под другим, справа скобка:]
Маккин
\Unclear{GELFAND-P000780}{bottom-added-name-makkin}{Маккин (чтение предположительно).}
\newline
Эро
\Unclear{GELFAND-P000780}{bottom-added-name-ero}{Эро (чтение предположительно).}

% Reproducible image-review audit: original 2408 x 3448 pixels; full frame viewed with detail original.
% Crop coordinates below are original-pixel xyxy, origin at upper left, right/bottom exclusive.
% All listed crops were enlarged with PIL.Image.Resampling.NEAREST and directly visually inspected.
% Review crop top-a: xyxy=(190, 35, 1050, 325); scale=3x.
% Review crop top-b: xyxy=(995, 35, 2015, 320); scale=3x.
% Review crop formula19a: xyxy=(90, 350, 1500, 490); scale=3x.
% Review crop formula19b: xyxy=(250, 530, 1410, 710); scale=3x.
% Review crop number19: xyxy=(100, 375, 206, 465); scale=10x.
% Review crop denominator-sign: xyxy=(950, 594, 1032, 650); scale=10x.
% Review crop middle-a: xyxy=(258, 728, 1606, 941); scale=3x.
% Review crop equations20-21: xyxy=(103, 962, 1402, 1208); scale=3x.
% Review crop middle-text-a: xyxy=(249, 1215, 1280, 1737); scale=3x.
% Review crop middle-text-b: xyxy=(1230, 1215, 2257, 1418); scale=3x.
% Review crop coincide: xyxy=(1585, 1224, 1894, 1338); scale=5x.
% Review crop backlund: xyxy=(1160, 1550, 1580, 1670); scale=5x.
% Review crop fourier-equations: xyxy=(440, 1765, 1430, 2080); scale=3x.
% Review crop lower-head: xyxy=(230, 2168, 1480, 2669); scale=3x.
% Review crop equation23: xyxy=(80, 2640, 1115, 2852); scale=3x.
% Review crop crossed23: xyxy=(1125, 2630, 1935, 2840); scale=5x.
% Review crop bottom: xyxy=(120, 2860, 1640, 3345); scale=3x.
% Review crop solution-cancelled-word: xyxy=(248, 2515, 493, 2680); scale=5x.
% Review crop names: xyxy=(423, 3145, 990, 3360); scale=5x.
% Review crop overwritten-u: xyxy=(491, 1817, 582, 1900); scale=10x.
% Review crop polynomial-glyph20: xyxy=(611, 990, 685, 1055); scale=10x.
% Review crop polynomial-glyph20-full: xyxy=(606, 964, 684, 1059); scale=10x.
% Review crop polynomial-glyph21: xyxy=(785, 1070, 856, 1144); scale=10x.
% Review crop fourier-symbol-top: xyxy=(1035, 1768, 1110, 1856); scale=10x.
% Review crop fourier-symbol-bottom: xyxy=(905, 1980, 966, 2050); scale=10x.
% Review crop where: xyxy=(342, 2462, 479, 2542); scale=5x.

\SourcePage{GELFAND-P000781}{6}
% Additional independent review: original image directly inspected.
% Enlargement resamples the original pixels; it does not add source detail.
% REVIEW independent matrix-left: box=[400, 2320, 1270, 2570], scale=3x, interpolation=nearest; original-pixel coordinates.
% REVIEW independent matrix-right: box=[1210, 2320, 2260, 2570], scale=3x, interpolation=nearest; original-pixel coordinates.
Другая форма:
\begin{equation*}
\Psi(x,t)=2^{-n}\sum_{m=0}^{n}c_m(t)
\underset{\substack{\uparrow\\\text{полиномы Эрмита}}}{H_{n-m}(x)}
\tag{25}
\end{equation*}
Тогда
\begin{equation*}
\ddot{c}_m(t)+mc_m(t)=0
\tag{26}
\end{equation*}

Нелинейное соответствие:
\begin{equation*}
\{x_j\}\leftrightarrow\{c_m\}
\end{equation*}
От корней к коэфф. [волнистое подчёркивание] — получ. лин. сист. дифф. ур-й.

Равновесные решения (24)
\begin{equation*}
\bar{x}_j=\sum_{k=1}^{n}{}'(\bar{x}_j-\bar{x}_k)^{-1}
\tag{27}
\end{equation*}
В этом случае $c_m\equiv0$,
[зачёркнута короткая математическая запись]
\Unclear{GELFAND-P000781}{equilibrium-cancelled-condition}{Возможно $\forall t$; зачёркнутые знаки неразборчивы.}
$m=1,2,\ldots,n$; $c_0=1$

полож. равновесия — нули полинома Эрмита порядка $n$

\underline{Следствие:}

нули полин. Эрмита удовл. (27).

(Stieltjes; 1885)

Без $\Psi_{tt}$ ур-е (22) — ур-е для полиномов Эрмита.

Малые осцилляции вокруг полож. равновесия:
[вся строка подчёркнута волнистой линией]
\begin{equation*}
x_j=\bar{x}_j+\varepsilon_j(t)
\tag{28}
\end{equation*}
\begin{equation*}
A_{jk}=\delta_{jk}\sum_{\ell=1}^{n}{}'
(\bar{x}_j-x_\ell)^{-2}
-(1-\delta_{jk})(\bar{x}_j-\bar{x}_k)^{-2}
\end{equation*}
\Unclear{GELFAND-P000781}{matrix-unbarred-xell}{В первом слагаемом записано $x_\ell$ без видимой черты.}
\begin{equation*}
\ddot{\varepsilon}_j(t)+A_{jk}\varepsilon_k(t)=0
\end{equation*}
\Unclear{GELFAND-P000781}{oscillation-missing-identity}{Записано $A_{jk}$ в приведённом уравнении.}

Соб. числа целые —
\begin{equation*}
Av^{(s)}=a_s v^{(s)},\qquad a_s=s-1,\quad s=1,2,\ldots,n.
\end{equation*}

\begin{center}
— · —
\end{center}

% Reproducible image-review audit: original 2408 x 3448 pixels; full frame viewed with detail original.
% Crop coordinates below are original-pixel xyxy, origin at upper left, right/bottom exclusive.
% All listed crops were enlarged with PIL.Image.Resampling.NEAREST and directly visually inspected.
% Review crop top: xyxy=(104, 44, 1420, 662); scale=3x.
% Review crop hermite-label: xyxy=(1100, 320, 1588, 466); scale=3x.
% Review crop correspondence-a: xyxy=(288, 678, 1130, 1090); scale=3x.
% Review crop correspondence-b: xyxy=(1080, 950, 2180, 1072); scale=3x.
% Review crop equilibrium: xyxy=(134, 1090, 2160, 1580); scale=3x.
% Review crop equilibrium-cancel: xyxy=(1180, 1384, 1344, 1475); scale=10x.
% Review crop stieltjes: xyxy=(283, 1575, 1520, 1905); scale=3x.
% Review crop hermite-oscillations: xyxy=(300, 1895, 2158, 2150); scale=3x.
% Review crop lower: xyxy=(130, 2165, 2208, 3020); scale=3x.
% Review crop matrix-first: xyxy=(800, 2335, 1420, 2555); scale=5x.
% Review crop eigenvalues: xyxy=(285, 2740, 1697, 2970); scale=3x.
% Review crop equation28-label: xyxy=(131, 2190, 266, 2295); scale=10x.

\SourcePage{GELFAND-P000782}{unlabelled}
% Additional independent review: original image directly inspected.
% Enlargement resamples the original pixels; it does not add source detail.
\underline{\underline{3/IV-78 г.}}\qquad\underline{Сем. Гельфанда.}

\underline{Д. Б. Фукс} (и \underline{И. М. Гельфанд})

\underline{Вычисление} \underline{когомологий} \underline{различных} \underline{алгебр Ли.}

$g$ — веществ. алгебра Ли;

$C_q(g)$ — группа $q$-мерных цепей
\begin{equation*}
C_q(g)\underset{\mathrm{def}}{=}\Lambda^q g
\qquad\text{(внешняя степень)}
\end{equation*}
\begin{equation*}
d:C_q(g)\to C_{q-1}(g)
\tag{1}
\end{equation*}
\begin{equation*}
\begin{aligned}
d(g_1\wedge\ldots\wedge g_q)
&=\sum_{1\leq s<t\leq q}(-1)^{s+t-1}
[g_s,g_t]\wedge g_1\wedge\ldots\\
&\qquad\wedge\hat{g}_s\wedge\ldots\wedge\hat{g}_t\wedge\ldots\wedge g_q
\end{aligned}
\tag*{(1)'}
\end{equation*}

% Reproducible image-review audit: original 2408 x 3448 pixels; full frame viewed with detail original.
% Crop coordinates below are original-pixel xyxy, origin at upper left, right/bottom exclusive.
% All listed crops were enlarged with PIL.Image.Resampling.NEAREST and directly visually inspected.
% Review crop heading: xyxy=(118, 2310, 1550, 2674); scale=3x.
% Review crop definitions: xyxy=(215, 2675, 1818, 3075); scale=3x.
% Review crop bottom-left: xyxy=(107, 3080, 1300, 3413); scale=3x.
% Review crop bottom-right: xyxy=(1220, 3190, 2408, 3413); scale=3x.
% Review crop exponent: xyxy=(1320, 3193, 1490, 3270); scale=8x.
% Review crop wedge-hats: xyxy=(1840, 3188, 2355, 3367); scale=5x.
% Source-clear heading date: 3/IV-78 г.; catalogue date is not substituted for visible handwriting.

\SourcePage{GELFAND-P000783}{11}
% Additional independent review: original image directly inspected.
% Enlargement resamples the original pixels; it does not add source detail.
% REVIEW independent dual-exponent: box=[385, 405, 770, 555], scale=5x, interpolation=nearest; original-pixel coordinates.
% Original PNG reviewed at native 2408x3448 pixels.
% Reproducible crops: original-coordinate half-open bbox (x0,y0,x1,y1), enlargement factor; no added source pixels.
% Review crop top-text: 160,20,2310,450; 3x; inspected.
% Review crop upper-math: 70,430,2150,895; 3x; inspected.
% Review crop middle-math: 90,875,2200,1805; 3x; inspected.
% Review crop fields: 90,1790,2260,2510; 3x; inspected.
% Review crop bottom: 100,2500,2390,3420; 3x; inspected.
% Review crop read-top-word: 420,27,675,125; 3x; inspected.
% Review crop dual-sentence: 305,225,1810,340; 3x; inspected.
% Review crop dual-parenthesis: 210,323,1550,440; 3x; inspected.
% Review crop eq02-exponent: 465,407,645,486; 10x; inspected.
% Review crop cancelled-no: 340,2885,575,2980; 5x; inspected.
% Review crop functionals: 730,555,1020,665; 5x; inspected.
% Review crop field-insertion: 954,1775,1250,1900; 5x; inspected.
% Review crop right-edge-set: 1945,2880,2408,3030; 5x; inspected.
% Review crop left-edge-upper: 0,0,105,1750; 3x; inspected.
% Review crop left-edge-lower: 0,1720,105,3448; 3x; inspected.
% Review crop gutter-fragment: 0,3090,85,3448; 10x; inspected.
Легко видно, что $d^2=0$ и отсюда возникают \underline{гомологии},
обозначаемые $H_q(g)$.

Двойственная конструкция даёт когомологии
(в непр. случае берутся непр.
\Unclear{GELFAND-P000783}{dual-parenthesis-abbreviation}{ф-лы / ф-ии})
\begin{equation*}
\delta:C^{q-1}(g)\to C^q(g) \tag{2}
\end{equation*}
$C^q(g)$ --- \Unclear{GELFAND-P000783}{cochain-abbreviation}{ф-лы / ф-ли} на $C_q(g)$

$H^q(g)$ --- гомологии двойств. компл. с дифф. $\delta$;

Скалярное произведение позволяет отождествить $C^q$ и $C_q$.
Получаем:
\begin{align*}
d &: C_q(g)\to C_{q-1}(g)\\
\delta &: C_q(g)\to C_{q+1}(g)
\end{align*}
Положим
\begin{equation*}
\Delta=d\delta+\delta d \tag{3}
\end{equation*}
(опер. Лапласа)

$\Delta$ перестановочен с $d$ и $\delta$.

Легко показать, что в каждом классе гомологий есть
ровно один \underline{гармонический} элемент.
\begin{equation*}
\Delta a=0\quad\Leftrightarrow\quad da=0\text{ и }\delta a=0 \tag{4}
\end{equation*}
Рассмотрим алгебру Ли [над строкой вставлено: (форм.)] вект. полей на $\mathbb{R}$ вида:
\begin{equation*}
f(x)\frac{d}{dx}, \tag{5}
\end{equation*}
где
\begin{equation*}
f(x)=\sum_{k\ge2}a_kx^k. \tag{6}
\end{equation*}
Положим
\begin{equation*}
e_n=x^{n+1}\frac{d}{dx}
\qquad\text{(элемент «размерности» }n\text{)}
\end{equation*}
Когомологии вычислены Л. Гончаровой, но с тех пор были ещё попытки
понять это.

Алгебра $W_1$ не имеет даже локальной группы Ли

[Под $W_1$ вписано $\{a(x)\frac{d}{dx}\}$; два коротких вертикальных штриха связывают надпись с $W_1$.]

[Зачёркнуто: \Unclear{GELFAND-P000783}{cancelled-word-before-group}{Но / Ново}.]
У группы формальных диффеом. алгебра Ли:
$\{a(x)\frac{d}{dx},\ a(0)=0\}$
\begin{equation*}
L_0\underset{\mathrm{def}}{=}\left\{a(x)\frac{d}{dx},\ a(0)=0\right\}
\quad\text{--- она алгебра Ли форм. группы} \tag{7}
\end{equation*}
\begin{equation*}
W_1\underset{\mathrm{def}}{=}\left\{a(x)\frac{d}{dx},\ a\text{ --- форм. ряд}\right\} \tag{8}
\end{equation*}
\begin{equation*}
L_1\underset{\mathrm{def}}{=}\left\{a(x)\frac{d}{dx},\ a(0)=0,\ a'(0)=0\right\} \tag{9}
\end{equation*}

[В нижней части слева, за границей листа:
\Unclear{GELFAND-P000783}{gutter-lower-fragment}{обрезанный фрагмент знака}.]

\SourcePage{GELFAND-P000784}{unlabelled}
% Additional independent review: original image directly inspected.
% Enlargement resamples the original pixels; it does not add source detail.
% Original PNG reviewed at native 2408x3448 pixels.
% Reproducible crops: original-coordinate half-open bbox (x0,y0,x1,y1), enlargement factor; no added source pixels.
% Review crop top: 115,25,2160,835; 3x; inspected.
% Review crop upper-middle: 115,815,2180,1410; 3x; inspected.
% Review crop formula: 95,1400,2315,1775; 3x; inspected.
% Review crop middle: 95,1760,2250,2440; 3x; inspected.
% Review crop bottom: 80,2420,2280,3195; 3x; inspected.
% Review crop reason-line: 224,350,1410,439; 3x; inspected.
% Review crop reference01: 110,395,390,477; 3x; inspected.
% Review crop definition-first: 125,548,1260,635; 3x; inspected.
% Review crop history-line: 265,740,1100,826; 3x; inspected.
% Review crop cancelled-reference02: 1910,855,2110,936; 5x; inspected.
% Review crop overalgebra: 1170,945,1650,1035; 3x; inspected.
% Review crop homology-trudnee: 214,1100,757,1195; 3x; inspected.
% Review crop speaker-heading: 255,1200,550,1315; 3x; inspected.
% Review crop generated-by: 440,1325,1078,1430; 3x; inspected.
% Review crop eq11-left: 435,1597,1062,1758; 10x; inspected.
% Review crop eq11-middle: 1030,1590,1667,1758; 10x; inspected.
% Review crop eq11-right: 1640,1602,2312,1795; 10x; inspected.
% Review crop weight-sentence: 260,1800,1865,1884; 3x; inspected.
% Review crop decomposition: 160,1890,690,2006; 3x; inspected.
% Review crop crossed-heading: 250,2750,1050,2870; 3x; inspected.
% Review crop overalgebra-full: 800,935,1550,1045; 3x; inspected.
% Review crop homology-first: 266,1020,2085,1136; 3x; inspected.
% Review crop word-before-elements: 375,1320,770,1428; 5x; inspected.
% Review crop weight-firstword: 260,1795,502,1885; 5x; inspected.
% Review crop definition-minimal: 223,625,1240,716; 3x; inspected.
% Review crop left-edge-upper: 0,0,105,1750; 3x; inspected.
% Review crop left-edge-lower: 0,1720,105,3448; 3x; inspected.
\begin{equation*}
L_0\subset W_1,\qquad \dim W_1/L_0=1.
\end{equation*}
$L_0$ интегрируется в группу

$W_1$ не интегрируется в группу.

Поэтому сущ. вся \Unclear{GELFAND-P000784}{reason-abbreviation}{следу.} история.

\underline{\Unclear{GELFAND-P000784}{reference-heading-one}{Спрв. / Справ.} 1.}

\underline{Алг. Ли} наз. алгеброй \underline{конечного дефекта}, если она
имеет подалг. конечной коразмерности, интегрир. в группу Ли.
Миним. коразм. наз. \underline{дефектом}.

\underline{С чего началось:} группы автоморфизмов.

\underline{\Unclear{GELFAND-P000784}{reference-heading-two}{Спрв. / Справ.} 2.}
Пусть нет центра; тогда группа автоморфизмов
[зачёркнуто: \Unclear{GELFAND-P000784}{cancelled-word-after-automorphisms}{имеет / неразборчивое слово}]
имеет алгебру Ли, явл.
\Unclear{GELFAND-P000784}{automorphism-algebra-word}{надалгеброй / подалгеброй} самой алгебры.

Когомологии $W_1$ и $L_0$ находятся легко (Гельфанд, Фукс);
а от $L_1$ труднее.

[Короткая подчёркнутая надпись с диагональной штриховкой:
\Unclear{GELFAND-P000784}{hatched-short-heading}{Фукс; / Факт;}.]

$L_1$ порождена элементами $e_1,e_2,\ldots$
\begin{equation*}
[e_n,e_m]=(m-n)e_{m+n} \tag{10}
\end{equation*}
\begin{equation*}
\begin{aligned}
d(e_{i_1}\wedge\cdots\wedge e_{i_q})
={}&\sum_{s<t}(-1)^{s+t-1}(i_t-i_s)e_{i_s+i_t}\wedge e_1\wedge\cdots\\
&\wedge\overset{'}{e}_{i_s}\wedge\cdots\wedge\overset{'}{e}_{i_t}\wedge\cdots\wedge e_{i_q}
\end{aligned} \tag{11}
\end{equation*}
\Unclear{GELFAND-P000784}{equation-eleven-first-factor}{После $e_{i_s+i_t}$ в источнике написано $e_1$.}
\Unclear{GELFAND-P000784}{equation-eleven-overmarks}{Над $e_{i_s}$ и $e_{i_t}$ стоят короткие штрихи; значение помет не установлено.}

Видно, что сумма индексов не меняется! Получается
\Unclear{GELFAND-P000784}{decomposition-word}{расслоение / разложение}
в прямую сумму конечномерных комплексов.

\underline{Теорема 1} (Л. Гончарова)

а) $\forall q\ge1\quad\Rightarrow\quad\underline{\dim H_q(L_1)=2}$

б) «вес» образующих
\begin{equation*}
\frac{3q^2\pm q}{2}
\end{equation*}
(сумма индексов)

\begin{center}
--- $\cdot$ ---
\end{center}

$C_{q,n}$ --- $q$-цепи веса $n$. Есть оператор Лапласа:
\begin{equation*}
\Delta:C_{q,n}\to C_{q,n} \tag{12}
\end{equation*}

[Подчёркнутая строка с диагональной штриховкой: Введем ещё структуру]
\Unclear{GELFAND-P000784}{hatched-structure-heading}{Функция диагональной штриховки неясна: выделение или зачёркивание.}

$\delta$ сопряженный по скал. произведению, в к-ром мономы
$\underset{i_1<\cdots<i_q}{e_{i_1}\wedge\cdots\wedge e_{i_q}}$
образуют ортонормированный базис.

\SourcePage{GELFAND-P000785}{12}
% Additional independent review: original image directly inspected.
% Enlargement resamples the original pixels; it does not add source detail.
% REVIEW independent eigenvalue-number: box=[320, 1535, 1330, 1740], scale=3x, interpolation=nearest; original-pixel coordinates.
% Original PNG reviewed at native 2408x3448 pixels.
% Reproducible crops: original-coordinate half-open bbox (x0,y0,x1,y1), enlargement factor; no added source pixels.
% Review crop top: 115,20,2310,850; 3x; inspected.
% Review crop upper-middle: 140,850,2250,1820; 3x; inspected.
% Review crop middle: 155,1800,2390,2400; 3x; inspected.
% Review crop bottom: 140,2350,2260,3430; 3x; inspected.
% Review crop third-basis: 730,142,967,238; 10x; inspected.
% Review crop margin-name: 177,365,478,484; 3x; inspected.
% Review crop cancelled-line: 310,852,1322,942; 3x; inspected.
% Review crop count-word: 1080,1454,1620,1545; 3x; inspected.
% Review crop multiplicities: 811,1712,1243,1820; 3x; inspected.
% Review crop sequences: 1490,1903,2095,2115; 5x; inspected.
% Review crop antisymmetric: 330,2910,1380,3020; 3x; inspected.
% Review crop sign-caption: 852,3110,1370,3215; 3x; inspected.
% Review crop date: 277,3310,1080,3427; 3x; inspected.
% Review crop count-claim: 1325,852,2320,940; 3x; inspected.
% Review crop cancelled-left: 320,857,761,931; 10x; inspected.
% Review crop cancelled-right: 721,848,1320,949; 10x; inspected.
% Review crop all-eigenvalues: 320,448,1220,550; 3x; inspected.
% Review crop fact-extra-mark: 374,345,485,440; 10x; inspected.
% Review crop top-firstline: 155,33,1115,131; 3x; inspected.
% Review crop combination: 1280,2900,1950,3030; 3x; inspected.
% Review crop left-edge-upper: 0,0,100,1750; 3x; inspected.
% Review crop left-edge-lower: 0,1720,100,3448; 3x; inspected.
3 кандидата на ортонорм. базис:
\begin{equation*}
e_k,\qquad\frac{e_k}{k!},\qquad\alpha_ke_k.
\end{equation*}
\Unclear{GELFAND-P000785}{third-basis-coefficient}{Коэффициент в третьем кандидате: $\alpha_k$ / $\lambda_k$.}
Здесь используется $e_k$

На внешние степени переносится.

Получаем опер. Лапласа (12).

\underline{Факты}
\Unclear{GELFAND-P000785}{facts-extra-mark}{После заголовка --- высокий заштрихованный знак, возможно III.}

А. \Unclear{GELFAND-P000785}{all-eigenvalues-mark}{Все / Все]} соб. значения --- целые!

Б. они перечисляются!

В пр-ве
\begin{equation*}
C_{*,n}=\bigoplus_q C_{q,n}
\end{equation*}

[Зачёркнуто: соб. числа
\Incomplete{GELFAND-P000785}{cancelled-spectral-phrase}{Средняя часть фразы неразборчива; справа возможно «способ…».}]

Число разных соб. значений $=$
\newline
$=$ числу разложений $n$ в сумму
\begin{equation*}
n=i_1+\cdots+i_q \tag{13}
\end{equation*}
\begin{equation*}
\begin{cases}
i_1\ge1\\
i_2-i_1\ge3\\
\cdots\\
i_q-i_{q-1}\ge3
\end{cases} \tag{14}
\end{equation*}
Разложению (13) отвечает \Unclear{GELFAND-P000785}{spectral-count-wording}{количество}
\begin{equation*}
C_{i_1}^3+\cdots+C_{i_q}^3-\sum_{s<t}i_si_t \tag{15}
\end{equation*}
соб. значений (кратность)

\underline{Упражнение:}

число (15) $\ge0$ и $=0$
\begin{equation*}
\Leftrightarrow\quad
\begin{aligned}
&\text{в }1,4,7,\ldots\qquad \Sigma=\frac{3q^2-q}{2}\\
&\text{или }2,5,8,\ldots\qquad \Sigma=\frac{3q^2+q}{2}
\end{aligned}
\end{equation*}
[После каждой из двух последовательностей, перед суммой справа, стоит короткая дуга.]
\Unclear{GELFAND-P000785}{exercise-arc-marks}{Значение двух дугообразных помет неясно.}

Соб. вектора выписываются

\underline{Символич. исчисление:}
\begin{equation*}
z=e_1,\quad e_2=z^2,\quad e_3=z^3;\ \ldots
\end{equation*}

\underline{Опис. соб. вектора:}
\begin{equation*}
\left(\prod_{1\le i<j\le q}(z_i-z_j)\right)^3
\end{equation*}
умножить на
\begin{equation*}
\prod z_i\qquad\text{или}\qquad\prod z_i^2
\end{equation*}
Получ. кососим. многочлен. Он предст. в виде комб.
\begin{equation*}
z_1^{i_1}\cdots z_q^{i_q}+\underbrace{\cdots}_{\text{с нужн. знаками}}
\end{equation*}
Заменяем на
\begin{equation*}
e_{i_1}\wedge\cdots\wedge e_{i_q}
\end{equation*}

\underline{След. заседание} 4/IV-78 г.



\end{document}
